\chapter{Observability, SRE, dan SLO}

\section*{Tujuan Pembelajaran}
\begin{itemize}
  \item Membedakan \textit{monitoring} dan \textit{observability} serta tiga sinyal telemetri: \textit{log}, metrik, dan \textit{trace}.
  \item Menetapkan SLI, SLO, dan \textit{error budget} untuk perjalanan pengguna yang kritis.
  \item Menjelaskan praktik SRE (\textit{Site Reliability Engineering}) yang menyeimbangkan kecepatan perubahan dan keandalan.
\end{itemize}

\section{Pendahuluan}

% TODO
% Cakupan: dasbor server Medikara Health selalu hijau, tetapi pasien tetap mengeluh gagal membuat janji temu; ukuran keandalan perlu berangkat dari pengalaman pengguna, bukan dari kondisi mesin.

\section{Monitoring dan Observability}

% TODO
% Cakupan: monitoring untuk masalah yang sudah diketahui dan observability untuk menjawab pertanyaan baru; tantangan menelusuri gangguan di sistem terdistribusi; peran observability bagi manajemen dan auditor.

\section{Log, Metrik, Trace, dan OpenTelemetry}

% TODO
% Cakupan: tiga sinyal telemetri dan kegunaan masing-masing; OpenTelemetry sebagai standar terbuka instrumentasi di bawah CNCF; biaya penyimpanan telemetri dan cara mengendalikannya.
% Pendalaman teknis: menginstrumentasi aplikasi contoh dengan OpenTelemetry dan menampilkan satu trace lintas dua layanan.

\section{SLI, SLO, dan SLA}

% TODO
% Cakupan: perbedaan SLI, SLO, dan SLA; cara memilih SLI yang mencerminkan pengalaman pengguna; perhitungan waktu tidak tersedia yang diizinkan untuk berbagai target; alasan target 100 persen hampir selalu keliru.
% Kasus Medikara: SLI untuk membuat janji temu, mengambil hasil laboratorium, dan mengirim klaim BPJS Kesehatan.

\section{Error Budget dan Kebijakannya}

% TODO
% Cakupan: error budget sebagai selisih antara 100 persen dan target SLO; kebijakan saat anggaran habis (misalnya menunda fitur baru); error budget sebagai alat negosiasi antara bisnis dan tim teknis.

\section{Praktik SRE: Toil, Alerting, Insiden, dan Postmortem}

% TODO
% Cakupan: SRE menurut buku Google SRE; toil dan batasnya; alerting yang bermakna berbasis SLO; penanganan insiden dan peran-perannya; postmortem tanpa menyalahkan; KPI dan pelaporan ke direksi disinggung singkat (dibahas rinci di IT40813).
% Kasus Medikara: postmortem gangguan Medikara+ saat program vaksinasi.

\section{Aktivitas: SLO Layanan Kritis dan Rancangan Observability}

% TODO
% Tujuan aktivitas: menetapkan SLI, SLO, dan kebijakan error budget untuk dua sampai tiga perjalanan pengguna kritis, beserta sinyal telemetri yang dibutuhkan.
% Peran dalam tim: anggota non-teknis menentukan arti layanan yang baik dari sisi pasien dan manajemen; anggota teknis menerjemahkannya menjadi SLI terukur, target SLO, dan kebutuhan telemetri; seluruh anggota menyepakati kebijakan error budget.
% Langkah: (1) pilih dua sampai tiga perjalanan pengguna kritis; (2) rumuskan SLI untuk tiap perjalanan; (3) tetapkan target SLO dan hitung error budget; (4) daftar sinyal telemetri yang dibutuhkan; (5) tulis kebijakan error budget satu halaman.
% Luaran dan tonggak capstone: C3, SLO layanan kritis dan rancangan observability.
% Berkas sources yang direncanakan: sources/chapter11/templat-slo.xlsx, sources/chapter11/templat-kebijakan-error-budget.md, sources/chapter11/contoh-opentelemetry/ (opsional)

\section{Ringkasan}

% TODO

\section{Latihan}

\begin{enumerate}
  \item Hitung waktu tidak tersedia yang diizinkan dalam satu bulan untuk SLO ketersediaan 99,5 persen dan 99,9 persen. Jelaskan dampak selisihnya bagi biaya.
  \item Mengapa target SLO 100 persen hampir selalu merupakan keputusan yang buruk?
  \item Usulkan satu SLI untuk layanan janji temu yang mencerminkan pengalaman pasien, bukan sekadar kondisi \textit{server}.
\end{enumerate}
